% !TEX root = ../main.tex
\corrigesChapitre{5}{Mécanismes réactionnels}
\corrigesSouples

\begin{blocLivre}{cor-5-1}
\begin{corrige}{1}{Lire un mécanisme}
\demarche{Repérer les entités formées puis consommées, puis additionner les actes élémentaires.}
\begin{questions}
\item Chaque acte met en jeu \res{\text{deux entités}} qui se rencontrent lors d'un choc.
\item \ce{NO3} est formé par l'acte 1, puis consommé par l'acte 2 : c'est \res{\text{l'intermédiaire réactionnel}}.
\item En additionnant : \ce{2 NO2 + NO3 + CO -> NO3 + NO + NO2 + CO2}, puis en simplifiant :
$\res{\ce{NO2 + CO -> NO + CO2}}$.
\item \res{\text{Non}} : l'acte 1 consomme deux molécules de \ce{NO2} et l'acte 2 n'en reforme qu'une. Le dioxyde d'azote
est globalement consommé : c'est un réactif, il figure d'ailleurs dans l'équation.
\end{questions}
\end{corrige}
\end{blocLivre}
\begin{blocLivre}{cor-5-2}
\begin{corrige}{2}{Repérer des charges partielles}
\demarche{Calculer $\Delta\chi$ pour chaque liaison ; l'atome le plus électronégatif porte la charge partielle $\delta^-$.}
\begin{questions}
\item Les atomes de chlore et de brome portent trois doublets non liants, l'atome d'oxygène deux et l'atome d'azote un.
\item C--Cl ($\Delta\chi=0{,}61$) : C$^{\delta+}$, Cl$^{\delta-}$ ; C--O (0,89) et O--H (1,24) : O$^{\delta-}$, C et H$^{\delta+}$ ;
H--Br (0,76) : H$^{\delta+}$, Br$^{\delta-}$ ; C--N (0,49) et N--H (0,84) : N$^{\delta-}$, C et H$^{\delta+}$.
\item Sites donneurs : les atomes porteurs de doublets non liants (Cl, O, Br, N). Sites accepteurs : l'atome de
carbone lié au chlore, l'atome d'hydrogène du groupe hydroxyle, celui du bromure d'hydrogène et ceux du groupe amino.
\item $\Delta\chi=2{,}55-2{,}20=0{,}35<0{,}4$ : la liaison C--H est \res{\text{considérée comme non polarisée}}.
\end{questions}
\schema[corrige]{05-ex02-cor}
\end{corrige}
\end{blocLivre}
\begin{blocLivre}{cor-5-3}
\begin{corrige}{3}{Lire des flèches courbes}
\demarche{Une flèche qui pointe vers un atome y crée une liaison ; une flèche qui part d'une liaison la rompt.}
\begin{questions}
\item La liaison \res{\text{N--H}} se forme ; la liaison \res{\text{H--Cl}} se rompt, et son doublet reste sur l'atome de chlore.
\item Site donneur : l'atome d'azote, porteur d'un doublet non liant. Site accepteur : l'atome d'hydrogène de
\ce{HCl}, qui porte une charge $\delta^+$ car $\chi(\ce{Cl})>\chi(\ce{H})$.
\item Il se forme l'ion ammonium \ce{NH4+} et l'ion chlorure \ce{Cl-}. Charge totale : $0+0=0$ avant l'acte et
$(+1)+(-1)=0$ après : \res{\text{la charge est conservée}}.
\item C'est une \res{\text{réaction acide-base}} : \ce{HCl} cède un ion \ce{H+} à \ce{NH3}. Couples \ce{HCl}/\ce{Cl-} et
\ce{NH4+}/\ce{NH3}.
\end{questions}
\end{corrige}
\end{blocLivre}
\begin{blocLivre}{cor-5-4}
\begin{corrige}{4}{Catalyseur ou intermédiaire ?}
\demarche{Un catalyseur est consommé puis reformé ; un intermédiaire est formé puis consommé.}
\begin{questions}
\item Étape 1 + 2 $\times$ étape 2 :
\[\ce{2 NO + O2 + 2 NO2 + 2 SO2 -> 2 NO2 + 2 NO + 2 SO3},\qquad\text{soit}\qquad\res{\ce{2 SO2 + O2 -> 2 SO3}}.\]
\item \ce{NO}, présent au départ, est consommé par l'étape 1 et reformé par l'étape 2 : c'est \res{\text{le catalyseur}}.
\ce{NO2} est formé par l'étape 1 et consommé par l'étape 2 : c'est \res{\text{l'intermédiaire réactionnel}}.
\item Catalyseur et réactifs forment une seule phase gazeuse : \res{\text{catalyse homogène}}.
\item La rencontre simultanée de trois molécules, avec une orientation et une énergie convenables, est très
improbable : l'étape 1 se décompose sans doute elle-même en actes élémentaires mettant en jeu deux entités.
\end{questions}
\end{corrige}
\end{blocLivre}
\begin{blocLivre}{cor-5-5}
\begin{corrige}{5}{Des chocs plus ou moins efficaces}
\demarche{Un choc est efficace s'il permet la formation des nouvelles liaisons : orientation et énergie doivent convenir.}
\begin{questions}
\item Pour former la liaison C--O, l'atome d'oxygène doit heurter l'atome de carbone, du côté opposé à l'atome de
chlore. À droite, l'ion heurte l'atome de chlore : \res{\text{l'orientation ne convient pas}}.
\item Le choc doit être \res{\text{suffisamment énergétique}} pour permettre la rupture de la liaison C--Cl pendant
que la liaison C--O se forme.
\item Deux fois plus d'ions par unité de volume : les chocs sont \res{\text{deux fois plus fréquents}} et la vitesse
de l'acte élémentaire double.
\item Les entités, plus lentes, se heurtent moins souvent et moins violemment : la proportion de chocs efficaces
diminue, la transformation \res{\text{ralentit}}.
\end{questions}
\end{corrige}
\end{blocLivre}
\begin{blocLivre}{cor-5-6}
\begin{corrige}{6}{Flèches courbes de réactions acide-base}
\demarche{Le site donneur est un doublet non liant de la base ; le site accepteur est l'atome d'hydrogène qui sera transféré.}
\begin{questions}
\item Acte 1 : \ce{CH3COOH}/\ce{CH3COO-} et \ce{H2O}/\ce{HO-}. Acte 2 : \ce{CH3NH3+}/\ce{CH3NH2} et \ce{H3O+}/\ce{H2O}.
\item Acte 1 : donneur, un doublet non liant de l'atome d'oxygène de \ce{HO-} ; accepteur, l'atome d'hydrogène du
groupe carboxyle ($\delta^+$, lié à un atome d'oxygène). Acte 2 : donneur, le doublet non liant de l'atome d'azote ;
accepteur, un atome d'hydrogène de l'ion oxonium.
\item Voir les schémas : dans chaque acte, une flèche crée la nouvelle liaison avec l'atome d'hydrogène et une
seconde flèche rompt l'ancienne liaison O--H.
\item Acte 1 : $0+(-1)=(-1)+0$ ; acte 2 : $0+(+1)=(+1)+0$ : \res{\text{la charge est conservée}}.
\end{questions}
\schema[corrige]{05-ex06}
\end{corrige}
\end{blocLivre}
\begin{blocLivre}{cor-5-7}
\begin{corrige}{7}{Un acte ou deux ?}
\demarche{Repérer la liaison polarisée C--Br : le carbone est accepteur, le brome garde le doublet de la liaison rompue.}
\begin{questions}
\item Donneur : un doublet non liant de l'atome d'oxygène de \ce{HO-}. Accepteur : l'atome de carbone, porteur d'une
charge $\delta^+$ car $\chi(\ce{Br})>\chi(\ce{C})$. Une flèche va du doublet vers le carbone, une autre de la liaison
C--Br vers le brome.
\item L'intermédiaire est \res{\text{le carbocation}} \ce{(CH3)3C+} : son atome de carbone central ne forme que trois
liaisons et porte une charge positive.
\item Acte 1 : une seule flèche, de la liaison C--Br vers l'atome de brome. Acte 2 : une flèche d'un doublet non liant
de l'ion hydroxyde vers l'atome de carbone chargé positivement.
\item $\res{\ce{(CH3)3C-Br + HO- -> (CH3)3C-OH + Br-}}$.
\end{questions}
\schema[corrige]{05-ex07}
\end{corrige}
\end{blocLivre}
\begin{blocLivre}{cor-5-8}
\begin{corrige}{8}{Addition sur une double liaison}
\demarche{Une double liaison contient un doublet liant supplémentaire, qui peut être mis en commun avec un site accepteur.}
\begin{questions}
\item La double liaison C=C comporte un second doublet liant, moins fortement lié, qui peut former une nouvelle
liaison : c'est \res{\text{un site donneur}}. Dans \ce{HBr}, $\chi(\ce{Br})>\chi(\ce{H})$ : l'atome d'hydrogène, $\delta^+$,
est \res{\text{le site accepteur}}.
\item Voir les schémas : acte 1, une flèche de la double liaison vers H et une de la liaison H--Br vers Br ;
acte 2, une flèche d'un doublet non liant de l'ion bromure vers le carbone chargé positivement.
\item L'intermédiaire est \res{\text{le carbocation}} \ce{CH3-CH^+-CH3}.
\item $\res{\ce{CH3-CH=CH2 + HBr -> CH3-CHBr-CH3}}$ : il se forme le \res{\text{2-bromopropane}}.
\end{questions}
\schema[corrige]{05-ex08}
\end{corrige}
\end{blocLivre}
\begin{blocLivre}{cor-5-9}
\begin{corrige}{9}{Fabriquer de l'éthanol}
\demarche{Additionner les trois actes ; l'entité consommée au début et reformée à la fin est le catalyseur.}
\begin{questions}
\item La somme des trois actes, après simplification, donne $\res{\ce{CH2=CH2 + H2O -> CH3-CH2-OH}}$.
\item Intermédiaires : \res{\ce{CH3-CH2+}} (carbocation) et \res{\ce{CH3-CH2-OH2+}}. L'ion oxonium, consommé par
l'acte 1 et reformé par l'acte 3, est \res{\text{le catalyseur}}.
\item Voir les schémas ci-dessous.
\item L'acte 3 transfère un ion \ce{H+} à l'eau : il forme l'éthanol et régénère le catalyseur. À 300 °C, les chocs
sont plus fréquents et beaucoup plus souvent efficaces : la transformation est assez rapide pour une production industrielle.
\end{questions}
\schema[corrige]{05-ex09-cor}
\end{corrige}
\end{blocLivre}
\begin{blocLivre}{cor-5-10}
\begin{corrige}{10}{Une catalyse par les ions du fer}
\demarche{Additionner les étapes, puis raisonner sur les charges des ions qui doivent se rencontrer.}
\begin{questions}
\item Étape 1 + étape 2 :
\[\ce{S2O8^2- + 2 Fe^2+ + 2 Fe^3+ + 2 I- -> 2 SO4^2- + 2 Fe^3+ + 2 Fe^2+ + I2},\]
soit \ce{S2O8^2- + 2 I- -> 2 SO4^2- + I2} : \res{\text{c'est bien l'équation de la réaction}}.
\item Les ions \ce{Fe^2+}, consommés par l'étape 1 et reformés par l'étape 2, sont \res{\text{le catalyseur}}. Introduits
au départ, les ions \ce{Fe^3+} réagissent d'abord selon l'étape 2, qui forme des ions \ce{Fe^2+} : le cycle démarre
simplement par l'autre étape.
\item Sans catalyseur, la réaction met en jeu deux anions, qui se repoussent : les chocs efficaces sont rares. Avec le
catalyseur, chaque étape fait réagir des ions de charges opposées, qui s'attirent : \res{\text{les chocs efficaces sont plus fréquents}}.
\item Le temps de demi-réaction est divisé par $\frac{40}{2}=20$ : la vitesse est multipliée par \res{\text{environ 20}}.
\end{questions}
\end{corrige}
\end{blocLivre}
\begin{blocLivre}{cor-5-11}
\begin{corrige}{11}{Chocs et température}
\demarche{Convertir les températures en kelvins, puis comparer l'effet sur la fréquence des chocs à l'effet observé sur la vitesse.}
\begin{questions}
\item $\sqrt{\dfrac{303{,}15}{293{,}15}}\approx\res{1{,}017}$ : la fréquence des chocs n'augmente que de 1,7 \%.
\item \res{\text{Non}} : la vitesse double alors que la fréquence des chocs n'augmente que de 2 \% environ. La
proportion de chocs efficaces est donc multipliée par $\frac{2}{1{,}017}\approx\res{2}$ : la température agit
surtout sur l'énergie des chocs.
\item La fréquence des chocs est proportionnelle au produit des deux concentrations : elle est multipliée par
$2\times2=\res{4}$, comme la vitesse de l'acte.
\item Chaque concentration est divisée par 5 : la fréquence des chocs est divisée par $5\times5=\res{25}$. Le
refroidissement diminue encore la proportion de chocs efficaces : la transformation est pratiquement arrêtée.
\end{questions}
\end{corrige}
\end{blocLivre}
\begin{blocLivre}{cor-5-12}
\begin{corrige}{12}{Le chlore et la couche d'ozone}
\demarche{Repérer l'entité régénérée à chaque cycle, puis passer de la masse de CFC à la quantité d'atomes de chlore.}
\begin{questions}
\item L'atome de chlore, consommé par l'acte 1 et reformé par l'acte 2, est \res{\text{le catalyseur}} ; \ce{ClO} est
\res{\text{l'intermédiaire réactionnel}}. Bilan : $\res{\ce{O3 + O -> 2 O2}}$.
\item Chaque atome de chlore est régénéré après avoir détruit une molécule d'ozone : il recommence le cycle un très
grand nombre de fois, tant qu'il n'est pas piégé.
\item Chaque molécule de \ce{CF2Cl2} contient deux atomes de chlore :
\[n(\ce{CF2Cl2})=\frac{\qty{1000}{g}}{\qty{120,9}{g.mol^{-1}}}=\qty{8,27}{mol}\qquad\text{et}\qquad n(\ce{Cl})=\qty{16,5}{mol},\]
\[n(\ce{O3})=\num{1.0e5}\times\qty{16,5}{mol}=\qty{1,65e6}{mol},\]
\[m(\ce{O3})=\qty{1,65e6}{mol}\times\qty{48,0}{g.mol^{-1}}\approx\res{\qty{7,9e7}{g}},\ \text{soit environ 80 tonnes}.\]
\item Les CFC sont très stables : ceux émis avant l'interdiction mettent des années à atteindre la stratosphère et y
persistent longtemps ; chaque atome de chlore libéré, catalyseur, continue d'agir jusqu'à son élimination.
\end{questions}
\schema[corrige]{05-ex12-cor}
\end{corrige}
\end{blocLivre}
\begin{blocLivre}{cor-5-13}
\begin{corrige}{13}{L'arôme de banane}
\demarche{Additionner les cinq actes, puis comparer la charge portée par l'atome de carbone avant et après l'acte 1.}
\begin{questions}
\item Bilan : $\res{\ce{CH3-COOH + R-OH -> CH3-COO-R + H2O}}$. L'ion \ce{H+}, consommé par l'acte 1 et régénéré par
l'acte 5, est \res{\text{le catalyseur}}. Les quatre autres espèces formées puis consommées sont les intermédiaires.
\item Donneur : un doublet non liant de l'atome d'oxygène de l'alcool ; accepteur : l'atome de carbone porteur de la
charge positive. La flèche va du doublet vers cet atome de carbone (schéma ci-dessous).
\item Dans l'acide, l'atome de carbone ne porte qu'une charge partielle $\delta^+$ : c'est un site accepteur peu
efficace. Après l'acte 1, il porte une charge positive entière : les chocs avec l'alcool sont \res{\text{bien plus souvent efficaces}}.
\item Le chauffage accélère la transformation (la température est un facteur cinétique). Le catalyseur ne modifie
pas l'état final ; pour cette réaction, presque sans effet thermique, la température ne le modifie pratiquement pas
non plus : on obtient \res{\text{la même quantité d'ester, plus vite}}.
\end{questions}
\schema[corrige]{05-ex13-cor}
\end{corrige}
\end{blocLivre}
\begin{blocLivre}{cor-5-14}
\begin{corrige}{14}{Le pot catalytique}
\demarche{Choisir les coefficients des étapes pour que les entités fixées sur le métal disparaissent du bilan.}
\begin{questions}
\item Le catalyseur est solide et les réactifs gazeux : \res{\text{catalyse hétérogène}}. Le nid d'abeilles offre une
très grande surface de contact entre les gaz et le métal dans un faible volume, donc de nombreux sites actifs.
\item On additionne 2 fois les étapes (1), (2), (3) et (4), et une fois l'étape (5) :
\[\mathrm{2\,CO+2\,NO}+6\,\ast\longrightarrow\mathrm{2\,CO_2+N_2}+6\,\ast,\quad\text{soit}\quad\res{\ce{2 CO + 2 NO -> 2 CO2 + N2}}.\]
Les intermédiaires sont $\mathrm{CO}^{\ast}$, $\mathrm{NO}^{\ast}$, $\mathrm{N}^{\ast}$ et $\mathrm{O}^{\ast}$. Les sites
libres, occupés puis libérés, sont \res{\text{régénérés}} : ils jouent le rôle de catalyseur.
\item Au-dessous de 250 à 300 °C, les réactions en surface sont trop lentes : les polluants traversent le pot sans
être transformés. Sur un trajet court, le pot reste froid pendant une grande partie du trajet.
\item Le plomb se fixe définitivement sur des sites, qui ne sont plus jamais libérés : le catalyseur est
\res{\text{empoisonné}} et il ne reste plus assez de sites libres pour fixer les réactifs.
\end{questions}
\end{corrige}
\end{blocLivre}
\begin{blocLivre}{cor-5-15}
\begin{corrige}{15}{La cinétique départage deux mécanismes}
\demarche{Appliquer les règles admises au seul acte lent, puis comparer avec l'effet de $[\ce{HO-}]$ sur le temps de demi-réaction.}
\begin{questions}
\item Mécanisme A : $\res{v=k\,[\mathrm{RBr}]\,[\ce{HO-}]}$. Mécanisme B : l'acte lent ne fait intervenir que RBr,
$\res{v=k\,[\mathrm{RBr}]}$.
\item 2-bromo-2-méthylpropane : $t_{1/2}$ ne dépend pas de $[\ce{HO-}]$, la vitesse non plus : \res{\text{mécanisme B}}.
Bromométhane : quand $[\ce{HO-}]$ double, $t_{1/2}$ est divisé par deux, la vitesse double : \res{\text{mécanisme A}}.
\item Loi d'ordre 1 : $k=\frac{\ln2}{t_{1/2}}=\frac{\ln2}{85}\approx\res{\qty{8,2e-3}{s^{-1}}}$.
\item Avec \ce{HO-} en large excès, $v=k'\,[\mathrm{RBr}]$ avec $k'=k\,[\ce{HO-}]$, donc
$t_{1/2}=\frac{\ln2}{k\,[\ce{HO-}]}$ est inversement proportionnel à $[\ce{HO-}]$ : pour \qty{0,050}{mol.L^{-1}},
deux fois moins que \qty{0,10}{mol.L^{-1}}, $\res{t_{1/2}=\qty{100}{min}}$.
\end{questions}
\end{corrige}
\end{blocLivre}
